% !TEX root = ../main.tex
% ---------------------------------------------------------------------
%  P1 -- Speicherhierarchie und Cache
% ---------------------------------------------------------------------
\vlsec{cP1}{P1}{Speicherhierarchie \& Cache}
\label{sec:cache}

\SW{Grundproblem} Speicher ist schnell \emph{oder} gross, nie beides
(Leitungslänge, Adressdecodierung, Preis/Bit). Jede Stufe:
ca.\ 3--10$\times$ langsamer, 10--100$\times$ grösser.

\begin{ztab}{@{}Xrr@{}}
Ebene (Raspberry Pi 5) & Grösse & Zugriff \\ \midrule
Register & $\sim$100\,B & 1\,Cycle \\
L1 (I und D getrennt) & 64\,kB/Core & 4 \\
L2 & 512\,kB/Core & 11 \\
L3 (shared) & 2\,MB & 35 \\
RAM & 4\,GB & $\sim$100 \\
SD-Karte (Swap) & 128\,GB & $\sim$100\,000 \\
\end{ztab}
Ohne Cache wäre eine 2.4-GHz-CPU faktisch eine 24-MHz-CPU.

\subsection{Locality of Reference}
\begin{description}
\item[Zeitlich] dieselbe Adresse bald wieder (Schleifenvariable, Stack).
\item[Räumlich] Nachbaradresse bald (linearer Code, Arrays).
\end{description}
\SW{Zahlen} 90--95\,\% der Befehle laufen linear weiter. Von den Sprüngen
$\sim$80\,\% kurze if-then-else, $\sim$20\,\% Schleifensprünge (davon
85\,\% genommen). \ra Der Befehlsstrom bleibt in einem kleinen
Adressfenster. Praxis: Hit Rate 66--99\,\%.

\subsection{Cache Line (typisch 64\,B)}
Drei Gründe, nie einzelne Bytes zu cachen:
\begin{enumerate}
\item \T{Latenz amortisieren}: RAM-Zugriff $\sim$100 Cycles, danach Burst fast gratis.
\item \T{Verwaltung}: Tag/Valid/Dirty pro 4\,B wären grösser als die Nutzdaten.
\item nutzt automatisch die räumliche Lokalität.
\end{enumerate}

\subsection{Adresszerlegung}
\label{sec:zerlegung}
\begin{center}\ttfamily\small
|\ \ Tag\ \ |\ Index\ |\ Offset\ |
\end{center}
\begin{description}
\item[Offset] (LSBs) Position \emph{innerhalb} der Line.
\item[Index] welches \emph{Set} -- reines Adressdecodieren, kein Suchen.
\item[Tag] Rest; mitgespeichert, sagt aus welchem RAM-Block die Line stammt.
\end{description}

\begin{rezept}[Rezept: Adresse zerlegen]
\label{rez:zerlegung}
\begin{enumerate}
\item $\text{Offset-Bits}=\log_2(\text{Linegrösse})$
\item $\#\text{Sets}=\dfrac{\text{Cachegrösse}}{\text{Line}\times\text{Assoziativität}}$,\quad
      $\text{Index-Bits}=\log_2(\#\text{Sets})$
\item $\text{Tag-Bits}=\text{Adressbreite}-\text{Index}-\text{Offset}$
\item Probe: $\#\text{Blöcke}\times\#\text{Sets}\times\text{Line}=$ Adressraum
\end{enumerate}
\end{rezept}

\begin{bsp}[Beispiel: L1 des Pi 5]
64\,kB, 64\,B Lines, 4-way.
$\#\text{Sets}=\frac{65536}{64\cdot4}=256$ \ra
6 Offset-Bit, 8 Index-Bit, Rest ($32-14=18$) Tag.
\end{bsp}

\SW{Zugriffsablauf} Index \ra Zeile ansteuern $\;\;$ Tag vergleichen
\emph{und} Valid prüfen $\;\;$ beides ok $=$ \T{Hit}, mit Offset das Wort
herausgreifen $\;\;$ sonst \T{Miss}: Line laden, Tag eintragen, Valid setzen.

\subsection{Mapping-Arten}
\begin{ztab}{@{}lXl@{}}
Art & Prinzip & Kosten \\ \midrule
Direct mapped & Index erzwingt genau \emph{eine} Line & 1 Komparator \\
Voll assoziativ & jeder Block in jede Line, kein Index & alle Tags, CAM \\
$n$-way set-ass. & Index wählt Set, $n$ freie Plätze darin & $n$ Komparatoren \\
\end{ztab}
\SW{CAM} Voll assoziativ braucht Content Addressable Memory: pro Bit
SRAM-Zelle (6\,T) $+$ EXOR-Komparator, Ausgänge auf ein Wired-NOR.
Schnellstmöglicher Vergleich, aber sehr teuer in Fläche und Strom
\ra heute nur noch für Kleinststrukturen (Micro-TLB).

\begin{falle}[Cache Trashing (nur Direct mapped)]
\label{bsp:trash}
Zwei aktive Codestellen mit \emph{gleichem Index}, aber verschiedenem Tag
werfen sich gegenseitig hinaus. Hit Rate $\to0$, der Cache bleibt zu
7/8 leer.\\
\T{Konstruieren:} Abstand $=k\cdot\frac{\text{Cachegrösse}}{\text{Assoziativität}}$.
Direct mapped, 8\,kB, 64\,B Lines \ra Index $=$ A12:A6; die Adressen
\R{0x0000} und \R{0x2000} kollidieren. Abhilfe: Assoziativität erhöhen.
\end{falle}

\subsection{Verwaltungsbits}
\begin{description}
\item[Valid] Nach Reset steht Müll im Cache, der zufällig Tag $0$ haben
  kann. „Unbenutzt“ $\neq$ „Tag $=0$“. Reset setzt Valid auf 0.
\item[Dirty] (nur D-Cache) Cache ist neuer als RAM. Nur dirty Lines müssen
  beim Verdrängen zurückgeschrieben werden.
\end{description}

\subsection{Replacement}
\begin{ztab}{@{}lX@{}}
FIFO & ältester Eintrag; ignoriert aktuelle Nutzung \\
LRU & am längsten nicht \emph{benutzt}; braucht Ordnungsinfo bei jedem Zugriff, ab 16-way teuer \\
Random & gratis, erstaunlich gut, \emph{immun} gegen pathologische Muster \\
NMRU & Random, aber die zuletzt benutzte Line ausgenommen \\
\end{ztab}
LRU und FIFO kann man mit einer ungünstigen Schleife gezielt zum Versagen
bringen, Random nicht.

\subsection{Write-Through vs.\ Write-Back}
\begin{ztab}{@{}lXX@{}}
 & Write-Through & Write-Back \\ \midrule
RAM & immer aktuell & zeitweise veraltet \\
Bandbreite & hoch (jeder Write) & minimal (nur Verdrängen) \\
Multicore/DMA & unproblematisch & braucht Kohärenzprotokoll \\
\end{ztab}
Schleife, die eine Variable 1000$\times$ hochzählt: Write-Through 1000
RAM-Zugriffe, Write-Back genau einer.

\subsection{Harvard}
Getrennte I- und D-Caches, damit Befehlsholen und Datenzugriff pro Cycle
\emph{parallel} laufen (sonst Stall). Der I-Cache braucht kein Dirty-Bit
und keine Write-Logik. Trennung nur auf L1, ab L2 \T{unified}
$=$ modifizierte Harvard-Architektur.

\subsection{Mittlere Zugriffszeit}
\begin{formel}[AMAT -- Average Memory Access Time]
\label{form:amat}
\[\text{AMAT}=t_{\text{Hit}}+\text{MissRate}\times\text{MissPenalty}\]
Mehrstufig von innen nach aussen verschachteln.
\end{formel}

\begin{bsp}[Beispiel: Pi 5, Hitraten 95/80/50\,\%]
\[4+0.05\,\bigl(11+0.20\,(35+0.50\cdot100)\bigr)=5.4\ \text{Cycles}\]
gegenüber 100 Cycles ohne Cache. Die Kurve kippt brutal, sobald L1 von
95 auf 90\,\% fällt.
\end{bsp}
